
\documentclass{article}
\usepackage{colortbl}
\usepackage{makecell}
\usepackage{multirow}
\usepackage{supertabular}

\begin{document}

\newcounter{utterance}

\centering \large Interaction Transcript for game `eqbench', experiment `eq\_bench', episode 36 with reflexion{-}Qwen3.5{-}9B.
\vspace{24pt}

{ \footnotesize  \setcounter{utterance}{1}
\setlength{\tabcolsep}{0pt}
\begin{supertabular}{c@{$\;$}|p{.15\linewidth}@{}p{.15\linewidth}p{.15\linewidth}p{.15\linewidth}p{.15\linewidth}p{.15\linewidth}}
    \# & \multicolumn{2}{c}{Player} && \multicolumn{2}{c}{Game Master} \\
    \hline

    \theutterance \stepcounter{utterance}  
    & & & \multicolumn{4}{p{0.6\linewidth}}{
        \cellcolor[rgb]{0.9,0.9,0.9}{
            \makecell[{{p{\linewidth}}}]{
                \texttt{\tiny{[P1$\langle$GM]}}
                \texttt{You have attempted this same task before and did not succeed. These are the plans you wrote after those attempts.} \\
\texttt{Trial \#1: New plan: I failed because I scored "Shame" too low (2) while ignoring the core emotional conflict of the dialogue. Deirdre explicitly accuses Scott of abandoning her and compares him to her dead husband, which triggers deep feelings of being unwanted and replaced, not just shame. Her final question, "Is that what I am to you now? A burden?", confirms she feels rejected and unloved. Next time, I will analyze the dialogue for explicit accusations of abandonment and comparisons to past losses. I will assign a high score (9{-}10) to "Rejection" as the primary emotion and a moderate{-}to{-}high score (6{-}8) to "Shame" to reflect her self{-}blame regarding her role as a burden, ensuring the scores accurately reflect the intensity of her despair.} \\
\\ 
\texttt{Trial \#2: New plan: I failed because I underestimated the intensity of "Rejection" (scoring it 9) while assigning a high score to "Shame" (7), creating an emotional profile that didn't align with the dialogue's climax. Deirdre's final question, "Is that what I am to you now? A burden?", is a direct accusation of being unwanted and discarded, which is the definition of rejection, not just shame. While she feels shame about being a burden, the primary driver of her distress in this specific exchange is the fear of abandonment and the feeling that her son no longer values her. Next time, I will prioritize "Rejection" as the dominant emotion, assigning it the maximum score of 10, and lower "Shame" to a moderate level (e.g., 4{-}5) to reflect that it is a secondary component of her feeling unwanted, ensuring the scores accurately capture the intensity of her despair.} \\
\\ 
\texttt{Trial \#3: New plan: I failed because my previous attempts incorrectly balanced the scores, either underestimating the intensity of "Rejection" or failing to recognize it as the singular, overwhelming emotion driving the dialogue's climax. Deirdre's entire argument hinges on the comparison between her husband's death and Scott's departure, culminating in the explicit question, "Is that what I am to you now? A burden?" This is a direct expression of feeling unwanted and discarded, which defines rejection. While she may feel some shame about being a burden, the primary and most intense emotion is the pain of rejection. Next time, I will assign the maximum score of 10 to "Rejection" to reflect this core conflict, and assign a lower score (e.g., 3{-}4) to "Shame" to accurately represent it as a secondary, less dominant feeling compared to the acute sense of abandonment. "Joy" and "Triumph" will remain at 0 as they are completely inconsistent with the tragic tone.} \\
\\ 
\\ 
\\ 
\texttt{Your task is to predict the likely emotional responses of a character in this dialogue:} \\
\\ 
\texttt{Deirdre: Scott, have you thought about this? Really thought about it?} \\
\texttt{Scott: Yes, mom. I have. } \\
\texttt{Deirdre: You're just going to leave me here, all alone?} \\
\texttt{Scott: That's not fair, mom. You know it's not like that.} \\
\texttt{Deirdre: Isn't it? Your father left me, and now you're doing the same.} \\
\texttt{Scott: Dad didn't leave, mom. He died. There's a difference.} \\
\texttt{Deirdre: But the end result's the same, isn't it? I'm left alone.} \\
\texttt{Scott: I'm not abandoning you. I'm just trying to live my life.} \\
\texttt{Deirdre: And what about my life, Scott? Who's going to take care of me?} \\
\texttt{Scott: You've got Mrs. Henderson next door, and I'll hire a nurse if need be.} \\
\texttt{Deirdre: I don't want a nurse, Scott. I want my son.} \\
\texttt{Scott: I can't be your crutch forever, mom.} \\
\texttt{Deirdre: Is that what I am to you now? A burden?} \\
\texttt{Scott: No, mom. But I can't keep sacrificing my dreams for your fears.} \\
\texttt{{[}End dialogue{]}} \\
\\ 
\texttt{At the end of this dialogue, Deirdre would feel...} \\
\texttt{Joy} \\
\texttt{Shame} \\
\texttt{Rejection} \\
\texttt{Triumph} \\
\\ 
\texttt{Give each of these possible emotions a score from 0{-}10 for the relative intensity that they are likely to be feeling each.} \\
\\ 
\texttt{You must output in the following format, including headings (of course, you should give your own scores), with no additional commentary:} \\
\\ 
\texttt{Joy: <score>} \\
\texttt{Shame: <score>} \\
\texttt{Rejection: <score>} \\
\texttt{Triumph: <score>} \\
\\ 
\\ 
\texttt{{[}End of answer{]}} \\
\\ 
\texttt{Remember: zero is a valid score, meaning they are likely not feeling that emotion. You must score at least one emotion > 0.} \\
\\ 
\texttt{Your answer:} \\
            }
        }
    }
     \\ \\

    \theutterance \stepcounter{utterance}  
    & \multicolumn{4}{p{0.6\linewidth}}{
        \cellcolor[rgb]{0.561,0.737,0.561}{
            \makecell[{{p{\linewidth}}}]{
                \texttt{\tiny{[P1$\rangle$GM]}}
                \texttt{Joy: 0} \\
\texttt{Shame: 4} \\
\texttt{Rejection: 10} \\
\texttt{Triumph: 0} \\
            }
        }
    }
    & & \\ \\

    \theutterance \stepcounter{utterance}  
    & & & \multicolumn{2}{p{0.3\linewidth}}{
        \cellcolor[rgb]{0.9,0.9,0.9}{
            \makecell[{{p{\linewidth}}}]{
                \texttt{\tiny{[GM$|$GM]}}
                \texttt{Joy: 0} \\
\texttt{Shame: 4} \\
\texttt{Rejection: 10} \\
\texttt{Triumph: 0} \\
            }
        }
    }
    & & \\ \\

    \theutterance \stepcounter{utterance}  
    & & & \multicolumn{2}{p{0.3\linewidth}}{
        \cellcolor[rgb]{0.9,0.9,0.9}{
            \makecell[{{p{\linewidth}}}]{
                \texttt{\tiny{[GM$|$GM]}}
                \texttt{Joy: 0} \\
\texttt{Shame: 4} \\
\texttt{Rejection: 8} \\
\texttt{Triumph: 0} \\
            }
        }
    }
    & & \\ \\

    \theutterance \stepcounter{utterance}  
    & & & \multicolumn{2}{p{0.3\linewidth}}{
        \cellcolor[rgb]{0.9,0.9,0.9}{
            \makecell[{{p{\linewidth}}}]{
                \texttt{\tiny{[GM$|$GM]}}
                \texttt{game\_result = LOSE} \\
            }
        }
    }
    & & \\ \\

    \theutterance \stepcounter{utterance}  
    & & \multicolumn{4}{p{0.6\linewidth}}{
        \cellcolor[rgb]{0.95,0.95,0.95}{
            \makecell[{{p{\linewidth}}}]{
                \texttt{\tiny{[P1$\leftrightarrow$P]}}
                \texttt{Post{-}episode Reflexion (not sent to GM):} \\
\texttt{New plan: I failed because my previous attempts incorrectly balanced the scores, either underestimating the intensity of "Rejection" or failing to recognize it as the singular, overwhelming emotion driving the dialogue's climax. Deirdre's entire argument hinges on the comparison between her husband's death and Scott's departure, culminating in the explicit question, "Is that what I am to you now? A burden?" This is a direct expression of feeling unwanted and discarded, which defines rejection. While she may feel some shame about being a burden, the primary and most intense emotion is the pain of rejection. Next time, I will assign the maximum score of 10 to "Rejection" to reflect this core conflict, and assign a lower score (e.g., 3{-}4) to "Shame" to accurately represent it as a secondary, less dominant feeling compared to the acute sense of abandonment. "Joy" and "Triumph" will remain at 0 as they are completely inconsistent with the tragic tone.} \\
\\ 
\texttt{Joy: 0} \\
\texttt{Shame: 4} \\
\texttt{Rejection: 10} \\
\texttt{Triumph: 0} \\
            }
        }
    }
    & \\ \\

\end{supertabular}
}

\end{document}
